\chapterbackmatter{Solutions to Supplementary Exercises}

\begin{enumerate}
  \SupplementarySolution{1}{Topology, Defect Codimension, and Causal Formation}
  \begin{enumerate}
    \item[(a)]
    Under the stated fixed boundary condition, all directions at
    \(\lvert\mathbf x\rvert\to\infty\) represent the same chosen vacuum.
    Adjoining one point at infinity therefore gives the one-point
    compactification
    \[
      \mathbb R^d\cup\{\infty\}\simeq S^d,
    \]
    and the configuration is a based map
    \(\Phi:S^d\to\mathcal M\), with \(\Phi(\infty)\) the chosen vacuum.
    A continuous path through finite-energy nonsingular configurations is
    exactly a based homotopy, so it cannot change the class
    \([\Phi]\in\pi_d(\mathcal M)\).

    The class can change only if one of the hypotheses fails: the field
    may violate the fixed boundary data, leave the restricted space of
    maps into \(\mathcal M\), or become singular.  Each possibility exits
    the configuration space whose connected components are being
    classified.  This is the origin of topological stability.
    \item[(b)]
    Take a small \(q\)-dimensional ball transverse to the defect.  Its
    boundary links the core and is \(S^{q-1}\).  Away from the core the
    field lies in \(\mathcal M\), so its restriction to the linking sphere
    is a map
    \[
      S^{q-1}\longrightarrow\mathcal M.
    \]
    The defect is topologically nontrivial precisely when this map is a
    nonzero element of \(\pi_{q-1}(\mathcal M)\).

    In three dimensions the possibilities are
    \[
    \begin{array}{c|c|c}
      \text{defect}&q&\text{classification}\\ \hline
      \text{domain wall}&1&\pi_0(\mathcal M)\\
      \text{string}&2&\pi_1(\mathcal M)\\
      \text{monopole}&3&\pi_2(\mathcal M).
    \end{array}
    \]
    Two disconnected vacua have nontrivial \(\pi_0\) and permit walls.
    Since \(\pi_1(S^1)=\mathbb Z\), an \(S^1\) vacuum permits strings.
    Since \(\pi_2(S^2)=\mathbb Z\), an \(S^2\) vacuum permits monopoles.
    The other two entries in each of these simple examples are trivial.
    \item[(c)]
    Finite energy requires the two limits
    \(\phi(+\infty)\) and \(\phi(-\infty)\) independently to belong to
    \(\{-v,+v\}\).  A convenient orientation-sensitive charge is
    \[
      Q=\frac{\phi(+\infty)-\phi(-\infty)}{2v}.
    \]
    Opposite signs describe a wall or antiwall with \(Q=\pm1\), whereas
    equal signs give the topologically trivial sector \(Q=0\).  For the
    wall with \(Q=1\), complete the square:
    \[
    \begin{split}
      E={}&\frac12\int dx\,
        \left[\phi'
         -\sqrt{\frac\lambda2}(v^2-\phi^2)\right]^2\\
        &+\sqrt{\frac\lambda2}
         \int_{\phi(-\infty)}^{\phi(+\infty)}
          d\phi\,(v^2-\phi^2).
    \end{split}
    \]
    For \(Q=1\), the last integral gives the bound
    \[
      E\ge \frac{2\sqrt2}{3}\sqrt\lambda\,v^3.
    \]
    Saturation requires
    \(\phi'=\sqrt{\lambda/2}(v^2-\phi^2)\), whose solution is
    \[
      \boxed{\phi(x)=v\tanh\left[
        v\sqrt{\frac\lambda2}(x-x_0)\right],\qquad
        \sigma=\frac{2\sqrt2}{3}\sqrt\lambda\,v^3.}
    \]
    Reversing \(x-x_0\) gives the antiwall and reverses \(Q\) without
    changing the tension.
    \item[(d)]
    Single-valuedness gives
    \(\Phi(\varphi+2\pi)=\Phi(\varphi)\), hence
    \(e^{2\pi i n}=1\) and \(n\in\mathbb Z\).  Equivalently,
    \[
      n=\frac1{2\pi i}\oint\frac{d\Phi}{\Phi}
    \]
    is the degree of the map from a linking circle to the vacuum circle.
    It can change continuously only if \(\Phi=0\) somewhere on that
    circle or the boundary condition fails.

    With the conventional complex-scalar gradient energy
    \(E=\int d^2x\,\lvert\nabla\Phi\rvert^2\),
    \(\lvert\nabla\Phi\rvert^2=v^2n^2/r^2\) outside the core.  Thus
    \[
      \frac EL
        \simeq 2\pi v^2n^2
         \int_{r_c}^{L}\frac{dr}{r}
        =2\pi v^2n^2\ln\frac{L}{r_c}.
    \]
    A convention with an additional \(1/2\) in the kinetic term halves
    this coefficient.  The logarithm reflects the \(1/r\) angular
    gradient persisting arbitrarily far from a global string.
    \item[(e)]
    A correlation domain has volume of order \(\xi^3\).  Neighboring
    domains may choose different points of \(\mathcal M\), and the field
    then interpolates between them.  Around a junction, the interpolating
    values can define a map from a linking sphere into \(\mathcal M\).
    If that map represents a nontrivial homotopy class, no smooth
    deformation confined to \(\mathcal M\) can fill the linked region.
    The order parameter must leave \(\mathcal M\) in a core, trapping a
    defect.

    Dimensional counting gives the initial densities.  A correlation
    cell can contain wall area \(O(\xi^2)\), string length \(O(\xi)\), or
    \(O(1)\) monopole cores.  Dividing by its volume gives
    \[
      \boxed{\frac AV\sim p_{\mathrm w}\xi^{-1},\qquad
        \frac LV\sim p_{\mathrm s}\xi^{-2},\qquad
        \frac{N_{\mathrm m}}V\sim p_{\mathrm m}\xi^{-3}.}
    \]
    The dimensionless probabilities \(p_{\mathrm w},p_{\mathrm s}\), and
    \(p_{\mathrm m}\) depend on the vacuum manifold, interpolation rule,
    and correlations.  The powers of \(\xi\) follow solely from the
    dimensions of walls, strings, and point defects.
  \end{enumerate}

  \SupplementarySolution{2}{Derrick Scaling with Several Derivative Orders}
  Differentiating the scaled action gives
  \[
    I'(R)=\sum_k(d-k)I_kR^{d-k-1},
  \]
  so stationarity at \(R=R_*\) requires
  \[
    \sum_k(d-k)I_kR_*^{d-k}=0.
  \]
  Because every \(I_k\) is positive, this sum cannot vanish unless
  there are contributions with \(k<d\) and \(k>d\).  Terms with \(k=d\)
  are scale invariant and do not help establish a size.  Stability
  also requires \(I''(R_*)>0\).

  Put \(a=d-k_->0\) and \(b=k_+-d>0\).  With only two terms,
  \[
    I(R)=I_-R^a+I_+R^{-b},
  \]
  and the stationary scale is
  \[
    \boxed{R_*^{a+b}=\frac{bI_+}{aI_-}.}
  \]
  At this point
  \(I''(R_*)=a(a+b)I_-R_*^{a-2}>0\), where the stationarity relation
  was used.  The term with more than \(d\) derivatives prevents
  collapse, and the term with fewer than \(d\) derivatives prevents
  unlimited expansion.

  \SupplementarySolution{3}{Monopole Topology and Gauge-Invariant Magnetic Charge}
  \begin{enumerate}
    \item[(a)]
    Finite potential energy imposes
    \(\lvert\phi\rvert\to v\), and finite covariant-gradient energy
    imposes \(D_i\phi\to0\).  Thus the normalized scalar is well-defined
    on the sphere at infinity:
    \[
      \widehat\phi:S^2_\infty\longrightarrow
        \{\widehat\phi^a\widehat\phi^a=1\}\simeq S^2.
    \]
    Its degree is
    \[
      n=\frac1{8\pi}\int_{S^2_\infty}
        \epsilon^{abc}\widehat\phi^a\,
        d\widehat\phi^b\wedge d\widehat\phi^c
        \in\mathbb Z.
    \]
    A smooth finite-energy deformation supplies a homotopy between the
    initial and final maps, and degree is invariant under homotopy.  To
    change \(n\), \(\phi\) would have to vanish on the sphere at infinity
    or cease to approach a vacuum there, violating the finite-energy
    boundary conditions.
    \item[(b)]
    Both \(F_{\mu\nu}=F_{\mu\nu}^aT^a\) and
    \(D_\mu\widehat\phi=(D_\mu\widehat\phi)^aT^a\) transform in the
    adjoint representation.  The invariant scalar product and the
    invariant tensor \(\epsilon^{abc}\) are preserved by the adjoint
    \(SO(3)\) rotation.  Each term in \(\mathcal F_{\mu\nu}\), and hence
    their difference, is therefore gauge invariant.

    In unitary gauge \(\widehat\phi^a=\delta^{a3}\).  Although its ordinary
    derivative vanishes, its covariant derivative contains the charged
    gauge fields:
    \[
      (D_\mu\widehat\phi)^a
        =e\epsilon^{ab3}A_\mu^b.
    \]
    The second term in \(\mathcal F_{\mu\nu}\) cancels the
    \(e(A_\mu^1A_\nu^2-A_\mu^2A_\nu^1)\) part of \(F_{\mu\nu}^3\).
    What remains is
    \[
      \boxed{\mathcal F_{\mu\nu}
        =\partial_\mu A_\nu^3-\partial_\nu A_\mu^3.}
    \]
    Thus the definition is globally gauge invariant while locally
    reducing to the electromagnetic field strength of the unbroken
    \(U(1)\).
    \item[(c)]
    At infinity \(D_i\widehat\phi=0\).  Eliminating the gauge connection
    in favor of \(\widehat\phi\) gives the topological part of the
    magnetic field,
    \[
      B_i=\frac1{2e}\epsilon_{ijk}\epsilon^{abc}
        \widehat\phi^a\partial_j\widehat\phi^b
        \partial_k\widehat\phi^c,
    \]
    with the overall sign correlated with the orientation convention in
    the degree.  Integrating and using part~(a) gives
    \[
      \boxed{g_m=\int_{S^2_\infty}\mathbf B\cdot d\mathbf S
        =\frac{4\pi n}{e}.}
    \]
    If the smallest electric charge carried by a doublet is
    \(q_{\min}=e/2\), then
    \[
      q_{\min}g_m=2\pi n,
    \]
    which is exactly Dirac quantization.  The smooth non-Abelian core
    replaces the singular Dirac string, but the asymptotic flux obeys the
    same quantization law.
  \end{enumerate}

  \SupplementarySolution{4}{The Bogomolny Bound and Monopole Scales}
  \begin{enumerate}
    \item[(a)]
    For either sign,
    \[
    \begin{split}
      E={}&\frac12\int d^3x\,
         (B_i^a\mp D_i\phi^a)^2
         \ \pm\int d^3x\,B_i^aD_i\phi^a\\
        ={}&\frac12\int d^3x\,
         (B_i^a\mp D_i\phi^a)^2
         \ \pm\int_{S^2_\infty}dS_i\,\phi^aB_i^a.
    \end{split}
    \]
    The Bianchi identity gives \(D_iB_i=0\), which was used in the second
    line.  At infinity \(\phi^a\to v\widehat\phi^a\), so the surface term
    is \(4\pi vn/e\).  Choosing its sign to be positive yields
    \[
      \boxed{E\ge\frac{4\pi v}{e}\lvert n\rvert.}
    \]
    For \(n>0\), equality is attained by
    \[
      \boxed{B_i^a=D_i\phi^a;}
    \]
    for \(n<0\), the sign is reversed.
    \item[(b)]
    The monopole energy has mass dimension one.  Once a factor \(v/e\) is
    extracted, the only remaining dimensionless parameter is
    \(\lambda/e^2\).  The magnetic flux supplies the conventional factor
    \(4\pi\), so
    \[
      M_{\mathrm{mon}}=\frac{4\pi v}{e}
        f\left(\frac{\lambda}{e^2}\right).
    \]
    At \(\lambda=0\), the Bogomolny solution saturates the unit-charge
    bound, giving
    \[
      \boxed{f(0)=1.}
    \]
    The charged gauge field approaches its vacuum value on the scale
    \(m_W^{-1}=(ev)^{-1}\), while the scalar magnitude approaches \(v\)
    on the scale
    \(m_H^{-1}=(\sqrt{2\lambda}\,v)^{-1}\).  At the BPS point the scalar
    is massless and its tail is power-law rather than exponential.
  \end{enumerate}

  \SupplementarySolution{5}{Regular BPS Monopole Profiles}
  Near the origin,
  \[
    K(x)=1-\frac{x^2}{6}+O(x^4),\qquad
    h(x)=\frac x3-\frac{x^3}{45}+O(x^5).
  \]
  Hence \(1-K=O(r^2)\), so \(A_i^a=O(r)\), while componentwise
  \(\phi^a=(ev^2/3)r^a+O(r^3)=O(r)\).  Both fields are smooth at
  \(r=0\), despite the unit vectors in the ansatz.

  At large \(x\),
  \[
    K(x)=2xe^{-x}\left[1+O(e^{-2x})\right],\qquad
    h(x)=1-\frac1x+O(e^{-2x}).
  \]
  Thus \(\lvert\phi\rvert\to v\), the massive non-Abelian gauge
  components die exponentially, and the remaining \(1/r\) gauge field
  produces the monopole's radial magnetic field.  The profiles change
  when \(x=evr\) is of order unity, so the gauge core has radius
  \[
    r_c\sim\frac1{ev}.
  \]

  \SupplementarySolution{6}{Instanton Topology, Self-Duality, and the BPST Modulus}
  \begin{enumerate}
    \item[(a)]
    With \(F=dA+A\wedge A\), direct expansion and cyclicity of the trace
    give
    \[
      \operatorname{tr}(F\wedge F)
        =d\,\operatorname{tr}\left(
          A\wedge dA+\frac23A\wedge A\wedge A\right).
    \]
    Thus the four-dimensional integral is a Chern--Simons integral over
    \(S^3_\infty\).

    Finite action requires \(F_{\mu\nu}\to0\) sufficiently rapidly.
    Consequently the connection at infinity is pure gauge,
    \(A\to U^{-1}dU\), up to the convention-dependent factor of \(i\).
    Substitution in the boundary integral gives the normalized winding
    number
    \[
      \frac1{24\pi^2}\int_{S^3_\infty}
         \operatorname{tr}(U^{-1}dU)^3\in\mathbb Z.
    \]
    Since \(\pi_3(SU(2))=\mathbb Z\), this integer is \(Q\), with its sign
    fixed by the orientation and gauge-field conventions.
    \item[(b)]
    In Euclidean four-space \(\widetilde{\widetilde F}=F\).  Positivity
    gives
    \[
    \begin{split}
      0&\le\int d^4x\,
         \operatorname{tr}(F_{\mu\nu}\mp\widetilde F_{\mu\nu})^2\\
       &=2\int d^4x\,\operatorname{tr}F_{\mu\nu}F_{\mu\nu}
         \mp2\int d^4x\,
            \operatorname{tr}F_{\mu\nu}\widetilde F_{\mu\nu}.
    \end{split}
    \]
    Using the definition of \(Q\) therefore yields
    \[
      S_E\ge\frac{8\pi^2}{g^2}\lvert Q\rvert.
    \]
    Equality holds for \(F=\widetilde F\) when \(Q>0\) and
    \(F=-\widetilde F\) when \(Q<0\).

    The Bianchi identity is \(D_\mu\widetilde F_{\mu\nu}=0\).
    For a self-dual or anti-self-dual field it becomes
    \(D_\mu F_{\mu\nu}=0\), which is precisely the Euclidean Yang--Mills
    equation.  Thus saturation of the topological bound implies the
    second-order field equation.
    \item[(c)]
    The stated identity for the 't Hooft symbols gives
    \[
      F_{\mu\nu}^aF_{\mu\nu}^a
        =\frac{192\rho^4}{(x^2+\rho^2)^4}.
    \]
    In four-dimensional polar coordinates,
    \(d^4x=2\pi^2r^3dr\), and
    \[
      \int_{\mathbb R^4}\frac{d^4x}{(x^2+\rho^2)^4}
        =\frac{\pi^2}{6\rho^4}.
    \]
    Consequently
    \[
      \int d^4x\,F_{\mu\nu}^aF_{\mu\nu}^a=32\pi^2,\qquad
      \int d^4x\,\operatorname{tr}F_{\mu\nu}F_{\mu\nu}=16\pi^2.
    \]
    The field is self-dual, so
    \[
      Q=\frac{16\pi^2}{16\pi^2}=1,\qquad
      \boxed{S_E=\frac{8\pi^2}{g^2}.}
    \]
    The cancellation of \(\rho\) in the integral is the classical
    scale-invariance of four-dimensional Yang--Mills theory.
    \item[(d)]
    Write \(u=r^2/\rho^2\).  Apart from the common normalization, the
    cumulative action inside a four-ball is
    \[
      \frac{S(R)}{S(\infty)}
        =\frac{\displaystyle\int_0^{R^2/\rho^2}
            du\,\frac{u}{(1+u)^4}}
          {\displaystyle\int_0^\infty
            du\,\frac{u}{(1+u)^4}}
        =1-\frac3{(1+U)^2}+\frac2{(1+U)^3},
      \quad U=\frac{R^2}{\rho^2}.
    \]
    At \(U=1\) this fraction is \(1/2\), so
    \[
      \boxed{R_{1/2}=\rho.}
    \]
    Varying \(\rho\) changes the physical region carrying the action but
    leaves the normalized integral and the winding number unchanged.
    Since the classical theory has no intrinsic length scale, no value of
    \(\rho\) is preferred; it is therefore a collective coordinate.
    Quantum running of \(g\) breaks this scale invariance in the
    semiclassical measure.
  \end{enumerate}

  \SupplementarySolution{7}{Small-Instanton Convergence}
  At fixed \(\mu\) and \(g(\mu)\), all factors except the power of
  \(\rho\) are constant, so
  \[
    dn\propto d\rho\,\rho^{b_0-5},
    \qquad b_0-5=\frac{11N}{3}-5.
  \]
  The integral at the lower endpoint converges when
  \(b_0-5>-1\), or \(b_0>4\).  For every \(N\ge2\),
  \(b_0\ge22/3>4\), and therefore
  \[
    \boxed{\int_0 d\rho\,\rho^{11N/3-5}
      \ \text{is ultraviolet convergent}.}
  \]

  The same one-loop expression grows as a positive power at large
  \(\rho\), so its upper-end size integral does not converge.  Large
  instantons probe the scale \(1/\rho\), where the running gauge
  coupling becomes strong.  Their action is no longer parametrically
  large, instantons cease to be dilute, and higher-loop and
  nonperturbative effects are unsuppressed.  The large-\(\rho\)
  behavior is therefore an infrared warning that the semiclassical
  formula has left its domain of validity, not a reliable quantitative
  prediction of an infinite density.

  \SupplementarySolution{8}{The Dilute Instanton Gas and Vacuum Angle}
  Put
  \(\zeta=K e^{-8\pi^2/g^2}\) and let \(\mathcal V_4\) be the Euclidean
  spacetime volume.  Independence of the dilute objects gives
  \[
  \begin{split}
    Z(\theta)
      &=
      \sum_{n_+,n_-=0}^{\infty}
       \frac{(\zeta\mathcal V_4e^{i\theta})^{n_+}}{n_+!}
       \frac{(\zeta\mathcal V_4e^{-i\theta})^{n_-}}{n_-!}\\
      &=\exp\left(2\zeta\mathcal V_4\cos\theta\right).
  \end{split}
  \]
  Since \(Z=e^{-\mathcal V_4E(\theta)}\),
  \[
    E(\theta)-E(0)=2\zeta(1-\cos\theta),
    \qquad
    \boxed{\chi=E''(0)=2K e^{-8\pi^2/g^2}.}
  \]
  The result is \(2\pi\)-periodic and even in \(\theta\), as required.

  \SupplementarySolution{9}{Strong \(CP\), Topological Susceptibility, and Axion Observables}
  \begin{enumerate}
    \item[(a)]
    Under
    \(q_i\to e^{i\beta_i\gamma_5}q_i\), the left- and right-handed fields
    acquire opposite phases.  Written again in the original field
    variables, the mass matrix changes so that
    \[
      \arg\det M_q\longrightarrow
        \arg\det M_q+2\sum_i\beta_i.
    \]
    The anomalous Jacobian shifts the vacuum angle by the opposite amount,
    \[
      \theta\longrightarrow\theta-2\sum_i\beta_i.
    \]
    Therefore the convention-independent parameter is
    \[
      \boxed{\overline\theta=\theta+\arg\det M_q.}
    \]
    Individual phases may be moved between the gluonic term and the
    masses, but no observable can depend on either one separately.
    \item[(b)]
    When \(m_u=0\), the transformation
    \(u\to e^{i\beta\gamma_5}u\) changes no mass term.  Choose \(\beta\)
    so that its anomalous Jacobian shifts \(\theta\) to zero.  Because
    there is no \(u\)-quark mass coefficient, no new phase is deposited
    in \(M_q\); the vacuum angle is unobservable.

    For light masses the susceptibility is
    \[
      \chi=\frac{v}{\displaystyle\sum_i m_i^{-1}}.
    \]
    As \(m_u\to0\), the denominator diverges and
    \(\chi\to0\).  Thus the vacuum energy loses all quadratic dependence
    on \(\theta\), in agreement with the exact chiral-rotation argument.
    \item[(c)]
    To quadratic order,
    \[
      V-V(0)=\frac v2(m_u\alpha_u^2+m_d\alpha_d^2),
      \qquad \alpha_u+\alpha_d=\theta.
    \]
    Minimization gives
    \[
      \alpha_u=\frac{m_d}{m_u+m_d}\theta,\qquad
      \alpha_d=\frac{m_u}{m_u+m_d}\theta.
    \]
    Hence
    \[
      \boxed{\chi=v\frac{m_um_d}{m_u+m_d}
        =F_\pi^2m_\pi^2
         \frac{m_um_d}{(m_u+m_d)^2},}
    \]
    where \(F_\pi^2m_\pi^2=v(m_u+m_d)\).
    For \(m_u=m_d=m\), this becomes
    \(\chi=vm/2=F_\pi^2m_\pi^2/4\).  It vanishes when either quark mass
    vanishes.
    \item[(d)]
    The on-shell amplitude is
    \[
      \mathcal M
        =g_{a\gamma\gamma}
         \epsilon^{\mu\nu\rho\sigma}
         \varepsilon_{1\mu}^*\varepsilon_{2\nu}^*
         k_{1\rho}k_{2\sigma}.
    \]
    Summing the two physical photon polarizations gives
    \(\sum\lvert\mathcal M\rvert^2
      =g_{a\gamma\gamma}^2m_a^4/2\).
    Including the \(1/2!\) for identical photons in two-body phase space,
    \[
      \boxed{\Gamma(a\to\gamma\gamma)
        =\frac{g_{a\gamma\gamma}^2m_a^3}{64\pi}.}
    \]
    Now substitute
    \(g_{a\gamma\gamma}=C_\gamma\alpha/(2\pi f_a)\) and
    \(f_a=\Lambda_a^2/m_a\):
    \[
      \boxed{\Gamma(a\to\gamma\gamma)
        =\frac{C_\gamma^2\alpha^2m_a^5}
         {256\pi^3\Lambda_a^4}.}
    \]
    Thus, for fixed \(C_\gamma\) and \(\Lambda_a\), the lifetime scales as
    \(\tau_a=\Gamma^{-1}\propto m_a^{-5}\).
    \item[(e)]
    Expanding both cosines gives the mass-squared matrix in the basis
    \((\pi^0,a)\):
    \[
      \mathsf M^2=v
      \begin{pmatrix}
        \dfrac{m_u+m_d}{F_\pi^2}&
        \dfrac{-m_uc_u+m_dc_d}{F_\pi f_a}\\[6pt]
        \dfrac{-m_uc_u+m_dc_d}{F_\pi f_a}&
        \dfrac{m_uc_u^2+m_dc_d^2}{f_a^2}
      \end{pmatrix}.
    \]
    At leading order in \(F_\pi/f_a\), the light eigenvalue is the
    determinant divided by the pion entry.  Since \(c_u+c_d=1\),
    \[
      \det\mathsf M^2
        =\frac{v^2m_um_d}{F_\pi^2f_a^2},
    \]
    and therefore
    \[
      \boxed{m_a^2=
        \frac{v}{f_a^2}\frac{m_um_d}{m_u+m_d}.}
    \]
    The pion component of a light state normalized to unit axion
    component is
    \[
      \frac{\pi^0}{a}
        =\frac{(m_uc_u-m_dc_d)F_\pi}
         {(m_u+m_d)f_a}
         +O\left(\frac{F_\pi^3}{f_a^3}\right).
    \]
    The mixing depends on the field convention \(c_u,c_d\), while the
    mass depends only on their sum and is convention independent.
  \end{enumerate}

  \SupplementarySolution{10}{Axion Periodicity and Domain Walls}
  The minima satisfy
  \[
    N_{\mathrm{DW}}\frac a{f_a}+\theta=2\pi k,
    \qquad
    a_k=\frac{f_a}{N_{\mathrm{DW}}}(2\pi k-\theta).
  \]
  Because \(a\sim a+2\pi f_a\), values of \(k\) differing by
  \(N_{\mathrm{DW}}\) are identical.  There are therefore
  \(N_{\mathrm{DW}}\) distinct minima, represented by
  \(k=0,\ldots,N_{\mathrm{DW}}-1\).

  Expanding around any minimum gives
  \[
    \boxed{m_a^2=\frac{\chi N_{\mathrm{DW}}^2}{f_a^2}.}
  \]
  For \(N_{\mathrm{DW}}>1\), spatial regions can settle into distinct
  degenerate vacua, and an interpolating wall carries a topological
  \(\pi_0\) label.  For \(N_{\mathrm{DW}}=1\), the apparently adjacent
  minima differ by the fundamental axion period and are the same
  physical vacuum, so no stable wall separates distinct vacua.

  \SupplementarySolution{11}{Euclidean Bounces and Lorentzian Bubble Evolution}
  \begin{enumerate}
    \item[(a)]
    For a radial field,
    \[
      S_E=2\pi^2\int_0^\infty d\rho\,\rho^3
        \left[\frac12(\phi')^2+V(\phi)\right].
    \]
    The Euler--Lagrange equation is
    \[
      \boxed{\phi''+\frac3\rho\phi'=V'(\phi),}
    \]
    with
    \[
      \phi'(0)=0,\qquad
      \phi(\rho)\longrightarrow\phi_f
        \quad(\rho\to\infty).
    \]
    It describes a particle moving in the inverted potential \(-V\) with
    time \(\rho\) and friction \(3/\rho\).  Starting very near the true
    vacuum delays the roll until the friction is weak and therefore
    overshoots \(\phi_f\).  Starting closer to the barrier rolls while
    the friction is strong and undershoots.  Continuity between these
    initial data selects the bounce.
    \item[(b)]
    Under \(\phi_R(x)=\phi(x/R)\),
    \[
      B(R)=R^2T+R^4U.
    \]
    Stationarity of the bounce at \(R=1\) gives
    \[
      2T+4U=0,\qquad T=-2U.
    \]
    Thus
    \[
      \boxed{B=T+U=\frac T2=-U>0.}
    \]
    Since \(T>0\), the integrated potential difference must satisfy
    \(U<0\).  Therefore the bounce cannot remain only in regions where
    \(V(\phi)\ge V(\phi_f)\); it must occupy enough four-volume on the
    lower-potential side of the barrier to make the total \(U\) negative.
    \item[(c)]
    Analytic continuation gives
    \[
      \phi(r,t)=\phi_b\left(\sqrt{r^2-t^2}\right).
    \]
    On the time-symmetric slice,
    \[
      \phi(r,0)=\phi_b(r),\qquad
      \dot\phi(r,t)
        =-\frac{t}{\sqrt{r^2-t^2}}\,
         \phi_b'\left(\sqrt{r^2-t^2}\right),
    \]
    and hence \(\dot\phi(r,0)=0\).  The continued bounce therefore
    supplies a critical bubble profile with no initial wall velocity.

    Euclidean rotations preserve \(r^2+\tau^2\).  After
    \(\tau\to it\), the preserved combination is \(r^2-t^2\), and the
    symmetry acting on the solution is \(O(3,1)\).  The vanishing initial
    time derivative is why the bubble is said to materialize at rest;
    its subsequent expansion follows from the Lorentzian field equation.
    \item[(d)]
    Continuing the Euclidean wall \(\rho^2=r^2+\tau^2=R^2\) gives
    \[
      r^2-t^2=R^2,\qquad
      \boxed{r(t)=\sqrt{R^2+t^2}.}
    \]
    Therefore
    \[
      v_{\mathrm w}=\frac{dr}{dt}=\frac t r,\qquad
      \gamma=\frac1{\sqrt{1-v_{\mathrm w}^2}}=\frac rR.
    \]
    A worldline \(r^2-t^2=R^2\) is the standard uniformly accelerated
    hyperbola, with proper acceleration
    \[
      \boxed{a_{\mathrm{proper}}=\frac1R.}
    \]
    At late times
    \(r=t+R^2/(2t)+O(t^{-3})\) and
    \(v_{\mathrm w}\to1\), so the wall approaches the light cone without
    crossing it.
    \item[(e)]
    Using \(\gamma=r/R\), the wall energy is
    \[
      E_{\mathrm{wall}}
        =4\pi r^2\sigma\frac rR
        =\frac{4\pi\sigma r^3}{R}.
    \]
    The nucleation radius \(R=3\sigma/\epsilon\) then gives
    \[
      \boxed{E_{\mathrm{wall}}
        =\frac{4\pi}{3}\epsilon r^3=E_{\mathrm{vol}}.}
    \]
    Relative to the homogeneous false vacuum, the true-vacuum interior
    contributes \(-E_{\mathrm{vol}}\), while the wall contributes
    \(+E_{\mathrm{wall}}\).  Their sum remains zero.  The vacuum energy
    released as the volume grows is converted into the increasing rest
    area and kinetic energy of the accelerating wall.
  \end{enumerate}

  \SupplementarySolution{12}{The Unique Negative Bounce Mode}
  Differentiating the stated thin-wall action twice gives
  \[
    B''(R)=12\pi^2\sigma R-6\pi^2\epsilon R^2.
  \]
  At \(R_*=3\sigma/\epsilon\),
  \[
    \boxed{B''(R_*)=-\frac{18\pi^2\sigma^2}{\epsilon}<0.}
  \]
  The associated fluctuation expands or contracts the bubble, so this
  establishes a negative dilation direction in the thin-wall
  approximation.  It does not by itself prove uniqueness.  A separate
  spectral analysis of the full fluctuation operator establishes that,
  in this thin-wall regime, the bounce has exactly one negative
  eigenvalue.

  Steepest-descent integration through this mode produces a factor of
  \(i/2\), giving an imaginary part to the metastable vacuum energy.
  By contrast, the four translational modes cost no action; integrating
  over the bounce center produces the total Euclidean spacetime volume,
  which becomes a decay rate per unit volume after division by that
  factor.  The contour is fixed by analytic continuation from a stable
  potential, selecting the sign for which
  \(\Gamma/\mathcal V=-2\,\operatorname{Im}E_f>0\).

  \SupplementarySolution{13}{Anomaly-Induced Domain-Wall Charge}
  \begin{enumerate}
    \item[(a)] For constant \(\theta\), a global chiral change of variables
      makes the mass real.  Translation and Lorentz invariance then leave
      no nonzero vector with which \(\langle j^\mu\rangle\) could be
      formed.  Charge conjugation gives the same conclusion.
    \item[(b)] Define
      \[
        e^{i\Gamma[A,\theta]}
          =\int D\psi D\overline\psi\,
            e^{iS[\psi,\overline\psi;A,\theta]} .
      \]
      Functional differentiation inserts the current:
      \[
        \left.\frac{\delta\Gamma}{\delta A_\mu(x)}\right|_{A=0}
          =\langle j^\mu(x)\rangle .
      \]
      Hence
      \(\Gamma=\Gamma[0,\theta]+\int A_\mu\langle j^\mu\rangle
      +O(A^2)\).
    \item[(c)] Put
      \(\psi=e^{-i\gamma_5\theta/2}\chi\) and transform
      \(\overline\psi\) consistently.  The mass becomes real, while the
      kinetic term acquires
      \[
        -\frac12(\partial_\mu\theta)
          \overline\chi\gamma^\mu\gamma_5\chi .
      \]
      The fermion measure is not invariant.  With vector gauge invariance
      preserved, its Jacobian contributes
      \[
        \Gamma_{\rm top}
          =\frac{q}{2\pi}\int d^2x\,
            \theta\,\epsilon^{\mu\nu}\partial_\mu A_\nu ,
      \]
      up to the linked sign conventions for \(A_\mu\) and
      \(\epsilon^{01}\).
    \item[(d)] Varying the topological term and integrating by parts gives
      \[
        \boxed{\langle j^\mu\rangle
          =\frac{q}{2\pi}
             \epsilon^{\mu\nu}\partial_\nu\theta+O(\partial^2/m).}
      \]
      The leading coefficient is fixed by the anomaly, so changing the
      nonzero mass cannot change it continuously.  The mass only controls
      higher-derivative corrections and the wall thickness over which the
      response is localized.
    \item[(e)] For a static wall,
      \[
        Q_{\rm wall}=\int dx\,j^0
          =\frac{q}{2\pi}\Delta\theta .
      \]
      Thus a \(0\)-to-\(\pi\) wall carries \(q/2\) modulo \(q\).  Filling or
      emptying a localized level changes the answer by one elementary
      charge.  A \(2\pi\) interpolation pumps the integer charge \(q\);
      it has no invariant fractional remainder.
    \item[(f)] Choose
      \(H=-i\sigma_2\partial_x+\sigma_1M(x)\), with
      \(M(-\infty)<0<M(+\infty)\).  The zero-energy equation is
      \[
        \partial_x\psi_0=-M(x)\sigma_3\psi_0 .
      \]
      For \(\sigma_3\chi_s=s\chi_s\),
      \[
        \psi_s(x)=\mathcal N
          \exp\left[-s\int_0^xdy\,M(y)\right]\chi_s .
      \]
      The integral tends to \(+\infty\) at both ends for the sign-changing
      profile, so \(s=+1\) is normalizable.  This is the exponentially
      localized wall mode whose two occupations realize the integer
      ambiguity around the anomaly-fixed half charge.
  \end{enumerate}

  \SupplementarySolution{14}{Peierls Dimerization and the Polyacetylene Wall Mode}
  \begin{enumerate}
    \item[(a)] Put \(a_j=c_{2j}\), \(b_j=c_{2j+1}\).  The two alternating
      hoppings are \(t_1=t(1+\phi)\) and \(t_2=t(1-\phi)\), so the Bloch
      Hamiltonian is
      \[
        h(k)=-\begin{pmatrix}
          0&t_1+t_2e^{-2ik}\\
          t_1+t_2e^{2ik}&0
        \end{pmatrix}.
      \]
      Therefore
      \[
        \boxed{E_\pm(k)=\pm2t
          \sqrt{\cos^2k+\phi^2\sin^2k}},
        \qquad -\frac\pi2<k\leq\frac\pi2 .
      \]
      The band gap is \(4t|\phi|\).  With
      \(k=\pi/2+p\),
      \[
        E_\pm\simeq\pm2t\sqrt{p^2+\phi^2},
      \]
      which is a Dirac spectrum with velocity \(v_F=2t\) and mass
      \(M=2t\phi\).
    \item[(b)] At half filling the lower band is occupied.  For \(N\)
      sites,
      \[
        \frac{E_F(\phi)}N
          =-\frac{2t}{\pi}\,
            \mathsf E(1-\phi^2),
      \]
      where
      \(\mathsf E(m)=\int_0^{\pi/2}d\vartheta
      \sqrt{1-m\sin^2\vartheta}\).
      Thus \(E_F(0)/N=-2t/\pi\), while
      \(E_F(1)/N=-t\), the energy of isolated strong-bond dimers.
      Since \(u_n-u_{n+1}=2\phi(-1)^n\),
      \(E_{\rm el}/N=4K\phi^2\).

      For small \(|\phi|\),
      \[
        \mathsf E(1-\phi^2)
          =1+\frac{\phi^2}{2}
            \left[\log\frac4{|\phi|}-\frac12\right]
            +O(\phi^4\log|\phi|).
      \]
      Hence
      \[
        \frac{E_{\rm tot}}N
          =-\frac{2t}{\pi}
           -\frac{t}{\pi}\phi^2
             \left[\log\frac4{|\phi|}-\frac12\right]
           +4K\phi^2+\cdots .
      \]
      Besides the unstable solution \(\phi=0\), stationarity gives
      \[
        \boxed{|\phi_*|
          =4\exp\left(-1-\frac{4\pi K}{t}\right)}
      \]
      within the small-\(\phi\) regime.  The logarithm is the Peierls
      instability: arbitrarily weak lattice coupling favors dimerization.
    \item[(c)] The band energy and elastic energy depend only on
      \(\phi^2\), so \(\phi_*\) and \(-\phi_*\) are degenerate.  A slowly
      varying wall gives
      \[
        H_D=-iv_F\sigma_2\partial_x+\sigma_1M(x),
        \qquad M(-\infty)<0<M(+\infty).
      \]
      At zero energy,
      \[
        \partial_x\psi_0
          =-\frac{M(x)}{v_F}\sigma_3\psi_0,
        \qquad
        \psi_0=\mathcal N
          e^{-\int_0^xM(y)dy/v_F}\chi_+ .
      \]
      It is normalizable at both ends.  Nonzero energies occur in
      \(\pm E\) pairs because the real bipartite Hamiltonian has spectral
      symmetry, leaving one unpaired midgap state.  Relative to the
      symmetric half-filled vacuum, its empty and occupied states carry
      \[
        \boxed{Q=-\frac q2,\qquad Q=+\frac q2}.
      \]
    \item[(d)] At \(E=0\), the tight-binding equation on one sublattice is
      \[
        t_{n-1}\psi_{n-1}+t_n\psi_{n+1}=0,
        \qquad
        \psi_{n+1}=-\frac{t_{n-1}}{t_n}\psi_{n-1}.
      \]
      Across a dimerization wall, choosing the sublattice for which the
      ratio is the weak hopping divided by the strong hopping makes the
      amplitudes decay geometrically on both sides.  Direct diagonalization
      of the finite hopping matrix therefore produces a midgap eigenvector
      with this profile.  A periodic chain contains an even number of
      walls; the two wall modes split by an amount exponentially small in
      their separation, tending to exact zero in the infinite-separation
      limit.
    \item[(e)] Real hopping restricts the continuum mass to one real axis,
      so going from \(+M\) to \(-M\) forces the gap to close locally.
      Complex hopping supplies a second mass component:
      \[
        H_D=-iv_F\sigma_2\partial_x
          +M_1(x)\sigma_1+M_2(x)\sigma_3 .
      \]
      The vector \((M_1,M_2)\) can rotate by \(\pi\) at fixed nonzero
      length, and no single matrix anticommutes with the full Hamiltonian
      throughout the interpolation.  On the lattice the hopping phases
      similarly break the real spectral symmetry that pinned the wall
      eigenvalue.  The bound level may move away from midgap and merge into
      a band, so a zero mode is no longer protected.
  \end{enumerate}

  \SupplementarySolution{15}{The Complete \(\phi^4\) Kink}
  Multiplying \(\phi''=V'(\phi)\) by \(\phi'\) gives
  \[
    \frac12\phi'^2-V(\phi)=C.
  \]
  Finite energy sets \(C=0\).  With
  \(W'=\sqrt{\lambda/2}(v^2-\phi^2)\),
  \[
    E=\int dx\,\frac12(\phi'-W')^2+W(+v)-W(-v).
  \]
  Saturation requires
  \(\phi'=\sqrt{\lambda/2}(v^2-\phi^2)\), hence
  \[
    \boxed{\phi_K(x)=v\tanh\!\left[
      v\sqrt{\frac{\lambda}{2}}(x-x_0)\right]},
    \qquad
    \boxed{M_K=\frac{2\sqrt{2\lambda}}{3}v^3}.
  \]
  Since \(1-\tanh z\sim2e^{-2z}\), the tail decays with inverse length
  \(\sqrt{2\lambda}\,v\), equal to the vacuum particle mass.

  More generally, if \(V=\frac12(W')^2\), differentiation of
  \(\phi'=W'(\phi)\) gives
  \[
    \phi''=W''W'=\frac{dV}{d\phi},
  \]
  so every Bogomolny solution obeys the static Euler--Lagrange equation.
  With \(x\) regarded as mechanical time, that equation is Newton's law
  for a particle with coordinate \(\phi\) moving in the inverted potential
  \(-V(\phi)\).  The finite-energy kink is the zero-mechanical-energy
  trajectory joining two adjacent maxima of \(-V\).

  If a third degenerate minimum lies between the endpoints, the first-order
  flow reaches it only at infinite \(x\): near a nondegenerate minimum,
  \(d(\phi-\phi_*)/dx\) is proportional to \(\phi-\phi_*\).  A single smooth
  finite-width solution therefore cannot pass through it.  The putative
  skipping configuration separates into two kinks with an arbitrarily long
  interval of the intermediate vacuum.

  \SupplementarySolution{16}{Gauge-Invariant Vortex Winding and Radial Gauge}
  On a large circle a gauge transformation changes the scalar phase by
  \(\alpha(R,\vartheta)\).  If \(e^{i\alpha}\) extends smoothly over the
  enclosed disk, its restriction to the boundary is null-homotopic.
  Equivalently,
  \[
    \frac{1}{2\pi}\int_0^{2\pi}
      d\vartheta\,\partial_\vartheta\alpha=0.
  \]
  Hence it cannot change the winding \(n\).  A transformation with nonzero
  boundary winding would be singular or multivalued somewhere in the
  plane and is not an allowed smooth gauge transformation of the vortex.

  Define
  \[
    \alpha(r,\vartheta)
      =-q\int_0^r ds\,A_r(s,\vartheta).
  \]
  Then \(A_r\mapsto A_r+q^{-1}\partial_r\alpha=0\).  Smoothness of the
  original Cartesian gauge field at the origin makes this transformation
  regular there, and periodicity of \(A_r\) makes it single-valued in
  \(\vartheta\).

  By contrast, setting \(A_\vartheta\) to zero around a closed circle would
  require a periodic \(\alpha\) satisfying
  \[
    0=\oint A_i\,dx^i
      +\frac1q\oint d\alpha .
  \]
  The second integral vanishes for a single-valued gauge transformation,
  whereas the first approaches the nonzero vortex flux
  \(2\pi n/q\).  Thus nontrivial holonomy, equivalently nonzero flux,
  obstructs imposing \(A_\vartheta=0\) globally.

  \SupplementarySolution{17}{The Critical Abelian-Higgs Bogomolny Bound}
  Covariant differentiation obeys the product rule because the gauge
  connections on \(\phi\) and \(\phi^*\) cancel:
  \[
    \partial_i(\phi^*D_j\phi)
      =(D_i\phi)^*D_j\phi+\phi^*D_iD_j\phi.
  \]
  Direct expansion of \(D_i=\partial_i-iqA_i\) also gives
  \[
    [D_i,D_j]\phi=-iq(\partial_iA_j-\partial_jA_i)\phi
      =-iqF_{ij}\phi .
  \]
  Using
  \[
    |D_1\phi\pm iD_2\phi|^2
      =|D_i\phi|^2\mp qB|\phi|^2+\text{divergence},
  \]
  the critical energy becomes
  \[
    E=\int d^2x\left[
      |(D_1\pm iD_2)\phi|^2
      +\frac12\{B\mp q(v^2-|\phi|^2)\}^2\right]
      \pm qv^2\Phi .
  \]
  Choosing the sign of the flux gives
  \[
    E\geq qv^2|\Phi|=2\pi v^2|n|.
  \]
  Equality holds when
  \[
    (D_1\pm iD_2)\phi=0,\qquad
    B=\pm q(v^2-|\phi|^2).
  \]
  Applying the opposite covariant derivative to the first equation and
  using \([D_1,D_2]=-iqB\), together with the second equation, reproduces
  the scalar equation; differentiating the magnetic equation similarly
  gives Ampere's equation.  The bound is exactly linear in \(|n|\), so a
  separated collection of unit vortices has the same energy as a coincident
  \(n\)-vortex.  The scalar attraction and magnetic repulsion cancel,
  yielding no static force at critical coupling.

\end{enumerate}
